\documentclass[11pt,a4paper]{article}
\usepackage[T1]{fontenc}
\usepackage{lmodern,microtype,mathtools,amssymb,amsthm,booktabs,tabularx,array,xcolor}
\usepackage[margin=26mm,top=24mm,bottom=24mm]{geometry}
\usepackage{fancyhdr,titlesec}
\usepackage[colorlinks=true,linkcolor=black!75!teal,urlcolor=black!65!teal,citecolor=black!70!teal]{hyperref}
\usepackage{bookmark}
\definecolor{ink}{HTML}{153B43}
\titleformat{\section}{\large\bfseries\color{ink}}{\thesection}{.65em}{}
\titleformat{\subsection}{\normalsize\bfseries\color{ink}}{\thesubsection}{.65em}{}
\newtheorem{theorem}{Theorem}[section]
\newtheorem{corollary}[theorem]{Corollary}
\newtheorem{proposition}[theorem]{Proposition}
\newcommand{\CT}{\operatorname{CT}}
\newcommand{\RhoG}{\mathcal R}
\pagestyle{fancy}\fancyhf{}\fancyhead[L]{\small\color{ink}Distance inertia and prime-power Motzkin differences}
\fancyhead[R]{\small\color{ink}Research manuscript}\fancyfoot[C]{\small\thepage}
\setlength{\headheight}{15pt}\setlength{\emergencystretch}{1.8em}\raggedbottom
\hypersetup{pdftitle={A counterexample to Graffiti 313 and exact prime-power differences of Motzkin numbers},pdfsubject={Distance inertia on triangle-free graphs and Lengyel's Conjecture 5.5}}
\begin{document}
\begin{center}
{\LARGE\bfseries\color{ink} A counterexample to Graffiti 313 and exact prime-power differences of Motzkin numbers\par}
\vspace{8pt}
{Research manuscript \enspace$\cdot$\enspace 15 November 2026}
\end{center}
\begin{abstract}
The Cartesian product of the Petersen graph and a square disproves
Graffiti conjecture 313, listed by Fajtlowicz under 27 August 1988. It has
40 vertices, 100 edges, independence number 14, and exactly seven negative
distance eigenvalues. For an
arbitrary connected graph $H$, multiplying by a crown graph makes the
ratio in the conjecture tend to $|V(H)|/b(H)$, where $b(H)$ is the
largest order of an induced bipartite subgraph. The ratio is unbounded
even on connected regular graphs and on explicit abelian Cayley graphs.
An independent arithmetic result settles the general-multiplier cases of
Lengyel's Conjecture~5.5. A binary coefficient identity and a ternary
expansion at a repeated root give the exact orders. On the ternary side,
the leading coefficient is nonzero for every fixed nonnegative offset,
giving an exact valuation formula for all multipliers prime to three.
\end{abstract}
\section{A Cartesian product obstruction for Graffiti 313}\label{sec:graffiti313}

Fajtlowicz listed Conjecture 313 under the triangle-free graphs heading dated
27 August 1988 \cite[p.~81]{fajt313}. For a finite connected graph $G$ with
at least one edge, write $e(G)=|E(G)|$, let $\alpha(G)$ be its independence
number, and let $\nu_-(D_G)$ be the number of strictly negative eigenvalues
of the ordinary graph distance matrix, counted with multiplicity. The assertion is
\begin{equation}\label{eq:313source}
             e(G)/\alpha(G)\le \nu_-(D_G) \qquad(G\text{ triangle-free}).
\end{equation}
Aouchiche and Hansen recorded it as Conjecture 7.14 and reported positive
experimental evidence \cite{ah313}. Their displayed $p_-(D)$ agrees with
the inertia notation in the original inequality. The word ``positive'' in their explanatory sentence is a
transcription error, resolved by Fajtlowicz's explicit definition in row 313.

For a graph $H$, let $b(H)$ be the maximum order of an induced bipartite
subgraph. Let $B_s$ be the crown graph $K_{s,s}$ with a perfect matching
deleted, $s\ge4$. Given $H$, the auxiliary question is to determine the
limiting ratio of the two sides of \eqref{eq:313source} on $H\square B_s$.

\begin{theorem}\label{thm:313limit}
Let $H$ be a connected graph on $v\ge1$ vertices with $e$ edges. Write
$J_v$ for the $v\times v$ all-ones matrix and $G_s=H\square B_s$. Then
\begin{align*}
 \alpha(G_s)&=s\,b(H), & e(G_s)&=s[2e+v(s-1)],\\
 \nu_-(D_{G_s})&=s-1+\nu_-\bigl(D_H+\tfrac32 J_v\bigr).
\end{align*}
In particular,
\[
 \lim_{s\to\infty}\frac{e(G_s)}{\alpha(G_s)\nu_-(D_{G_s})}
 =\frac{v}{b(H)}.
\]
Every connected, triangle-free, nonbipartite $H$ gives strict violations
of \eqref{eq:313source} for all sufficiently large $s$. Moreover, the
ratio in this limit formula is unbounded over all connected triangle-free
graphs.
\end{theorem}

\begin{proof}
One-coordinate edges in a Cartesian product give the stated edge count.
The independence identity is the crown case of Hagauer and Klav\v zar's bipartite product
formula \cite[Prop.~7.3]{ikr313}. Its short proof identifies the role of
$b(H)$. Choose a perfect matching of $B_s$, for example by pairing the
left vertex $i$ with right vertex $i+1$ modulo $s$. For any independent
set in $G_s$, its intersections with the two $H$-fibres of each matched
edge project to disjoint independent sets in $H$. They contain together
at most $b(H)$ vertices. Conversely, alternate the two color classes of a
maximum induced bipartite subgraph of $H$ between the two sides of $B_s$.
The resulting independent set has exactly $s\,b(H)$ vertices.

The distance matrix of a Cartesian product satisfies the standard identity
\[
 D_{G_s}=D_H\otimes J_{2s}+J_v\otimes D_{B_s},
\]
since the distance is the sum of the factor distances. For prior spectral formulas see Indulal and the treatment of product matrices by Catral et al. \cite{indulal313,catral}. Here $H$
may have nonconstant distance row sums. In the natural order of the two
sides of the crown, $D_{B_s}$ has diagonal blocks $2(J_s-I_s)$ and
off-diagonal blocks $J_s+2I_s$. On the all-ones vector it has eigenvalue
$3s$. On the antisymmetric constant vector it has eigenvalue $s-4$.
For $x\perp\mathbf1_s$, it has eigenvalue zero on $(x,x)$ and eigenvalue
$-4$ on $(x,-x)$. For $s\ge4$ only the last type is negative. Thus on
$\mathbb R^v\otimes\langle\mathbf1_{2s}\rangle$ the product distance matrix
acts as $2s(D_H+\tfrac32J_v)$. For each nonconstant crown eigenvector
with eigenvalue $\lambda$, the corresponding block is $\lambda J_v$,
which adds one negative eigenvalue exactly if $\lambda<0$. This gives the
claimed inertia formula for $s\ge4$.

Taking the ratio gives the limit. Every triangle in a Cartesian product lies
in a single factor fibre. If $H$ is triangle-free and nonbipartite, $G_s$
is therefore triangle-free and $b(H)<v$. In fact with
$c=\nu_-(D_H+\tfrac32J_v)$ the exact criterion is
\[
             (v-b(H))(s-1)>b(H)c-2e.
\]
Finally, the triangle-free graphs of Erd\H{o}s \cite{erd313} have arbitrarily
small $\alpha(H)/v$. If necessary, connect their components by bridges,
which create no triangle and cannot increase the independence number.
Since $b(H)\le2\alpha(H)$, the possible limits $v/b(H)$ are unbounded.
\end{proof}

\begin{corollary}\label{cor:313eighty}
The connected triangle-free graph $C_5\square B_8$ has 80 vertices,
360 edges, independence number 32 and eleven negative distance eigenvalues.
In particular its two sides in \eqref{eq:313source} are $45/4$ and $11$.
For $C_5\square B_s$ with $s\ge4$, the excess of the left side over the
right side is $(s-7)/4$.
\end{corollary}
\begin{proof}
Here $b(C_5)=4$, since removing one vertex leaves a path. Write $A$ for
the adjacency matrix of $C_5$. Then $D_{C_5}=2J_5-2I_5-A$. On the orthogonal
complement of the constants the four eigenvalues of $D_{C_5}$ are
$-2-2\cos(2\pi j/5)$ for $1\le j\le4$, all strictly negative. Its row sums
are six, hence $\nu_-(D_{C_5}+\tfrac32J_5)=4$. The theorem gives
$e/\alpha=5(s+1)/4$ and $\nu_-=s+3$.

At $s=8$ the same decomposition gives the full characteristic polynomial
\[
 t^{67}(t-216)(t-20)(t+20)^7(t^2+48t+256)^2.
\]
The two roots of the quadratic arise from the two pairs of nonconstant
Fourier modes of $C_5$. Thus the inertia, in positive, negative and zero
order, is $(2,11,67)$.
\end{proof}
\begin{proposition}[A 40-vertex counterexample]\label{prop:petersen}
The Cartesian product of the Petersen graph and $C_4$ has 40 vertices,
100 edges, independence number 14 and seven negative distance eigenvalues.
\end{proposition}
\begin{proof}
Write $P=KG(5,2)$ for the Petersen graph. Its vertices are the ten
two-element subsets of a five-element set, with adjacency given by
disjointness. A triangle would use six elements. Two distinct intersecting
pairs have a common neighbour consisting of the two unused elements,
so the diameter is two. For the adjacency matrix $A$, counting common
neighbours gives $A^2=2I-A+J$. The row sums and the trace give its
spectrum $3,1^{(5)},(-2)^{(4)}$. Thus $D_P=2J-2I-A$ has spectrum
$15,(-3)^{(5)},0^{(4)}$. The distance spectrum of $C_4$ is
$4,(-2)^{(2)},0$. Each factor has constant distance row sum. The
constant product mode gives $4\cdot15+10\cdot4=100$.
A nonconstant mode of either factor with a constant in the other
gets multiplied by the order of the other factor. Both all-ones
matrices kill a doubly nonconstant mode. The product's distance
characteristic polynomial is therefore
\[
 t^{32}(t-100)(t+12)^5(t+20)^2.
\]

A pairwise-intersecting set of edges of $K_5$ lies in a star or a
triangle. Indeed two intersecting edges can be written $xy,xz$.
Every other edge contains $x$ or is $yz$. If $yz$ is present, the
three triangle edges exclude any further edge. Consequently every
independent set of size four in $P$ is a full star, and any two such
stars intersect. Thus $b(P)\le7$, and equality holds by taking the full
star at $1$ together with the mutually intersecting edges $23,24,25$.
The matching proof in Theorem~\ref{thm:313limit} applies to $C_4$,
so $\alpha(P\square C_4)=2b(P)=14$. Finally the product is
5-regular on 40 vertices and $100>14\cdot7$.
\end{proof}
\section{How large the violation can be}\label{sec:ram}
Write $\RhoG(G)=e(G)/(\alpha(G)\nu_-(D_G))$ for a connected graph of order
at least two. Its denominator is positive since a nonsingular principal
two-by-two distance submatrix has one negative eigenvalue.
The preceding construction quantifies the failure even with symmetry constraints.
\begin{theorem}\label{thm:growth}
There is an absolute constant $c>0$ and an infinite sequence of connected
regular triangle-free graphs of orders $N\to\infty$ satisfying
\[
 \RhoG(G)\ge c\,N^{1/4}/\sqrt{\log N}.
\]
There is also an explicit sequence of connected triangle-free Cayley
graphs on elementary abelian $2$-groups satisfying $\RhoG(G)\ge cN^{1/6}$.
\end{theorem}
\begin{proof}
For a connected base of order $v\ge4$, choose $s=v$ in
Theorem~\ref{thm:313limit}. Then $N=2v^2$ and
$\nu_-(D_H+\tfrac32J_v)\le v-1$ by testing on the all-ones vector.
Since $b(H)\le2\alpha(H)$, the exact product formulas give
\[
 \RhoG(H\square B_v)
 \ge \frac{v(v-1)}{b(H)(2v-2)}
 =\frac{v}{2b(H)}\ge\frac{v}{4\alpha(H)}.
\]
Alon, Ben-Shimon and Krivelevich \cite[Thm.~1.1]{abs}
construct regular triangle-free graphs of arbitrarily large order $t$
with independence number at most $C\sqrt{t\log t}$. For each such
graph choose a connected component $H$ with
$\alpha(H)/v\le C\sqrt{(\log t)/t}$. Such a component exists by taking
a weighted average of the component independence ratios. Its order $v$
tends to infinity since $\alpha(H)\ge1$. For $3\le v\le t$ the function
$(\log x)/x$ is decreasing, so $\alpha(H)\le C\sqrt{v\log v}$.
The product with $B_v$ is connected, regular and triangle-free. The
displayed estimate proves the first assertion.

For the explicit construction use Alon's binary Cayley graphs
\cite[Thm.~2.1]{alon}. They have $v=2^{3k}$ vertices for $3\nmid k$,
degree $d\asymp v^{2/3}$, and all nontrivial adjacency eigenvalues have
absolute value at most $L=O(v^{1/3})$. They are connected for all
sufficiently large $k$ by $L<d$. For an independent set $S$ of size $u$,
decompose its indicator into its constant part and its orthogonal part.
Its adjacency quadratic form is zero, whence
$0\ge d u^2/v-L(u-u^2/v)$ and $u\le vL/(d+L)=O(v^{2/3})$.
The bound above yields $\RhoG(G)\gg v^{1/3}\asymp N^{1/6}$.
On $\mathbb F_2\times\mathbb F_2^{3k}$ the crown $B_v$ has generators
$(1,x)$, $x\ne0$. The Cartesian product is therefore itself a Cayley
graph on $\mathbb F_2^{6k+1}$.
\end{proof}
\section{Lengyel's prime-power conjecture}

Let $M_m$ count paths of length $m$ from zero to zero, with steps $-1,0,1$, staying nonnegative. Write $P(x)=x^{-1}+1+x$ and $c_{m,j}=[x^j]P(x)^m$.
Reflection of the paths at their first visit to $-1$ gives
$M_m=c_{m,0}-c_{m,2}$. Equivalently,
$M_m=\operatorname{CT}((1-x^2)P(x)^m)$, as follows also
by expanding the Catalan sum for $M_m$.
Denote the $p$-adic valuation by $v_p$. The following theorem gives
the two general-multiplier assertions of Lengyel's Conjecture~5.5.
The source is \cite[Conj.~5.5]{len13}. Its third assertion for primes $p\ge5$ and multiplier $1$ is
Lengyel's Theorem~5.7 in \cite{len13}, which gives $v_p(M_{p^{n+1}}-M_{p^n})=n$ for every $p\ge5$ and $n\ge2$.

\begin{theorem}\label{thm:lengyel55}
For each positive integer $a\equiv1\pmod4$ and each odd $n\ge3$,
\[
 v_2(M_{a2^{n+1}}-M_{a2^n})=n .
\]
For each positive integer $a$ prime to $3$, put
$V=v_3\binom{2a}{a}$. Whenever $n\ge1$ and
$(3^n+1)/2\ge 2n+V$,
\[
 v_3(M_{a3^{n+1}}-M_{a3^n})=n+V .
\]
\end{theorem}

\subsection{A dyadic coefficient identity}

Put $m=a2^n$, $h=2^{n-1}$, and define
\[
 L_s=c_{s,1}-c_{s,3},\qquad X_m=L_{m-1}-2M_{m-1}.
\]
The identity $P(x)^2=P(x^2)+2P(x)$ gives the exact formula
\[
 M_{2m}-M_m=c_{m,2}-c_{m,1}
 +\sum_{k=1}^m 2^k\binom{m}{k}
 \operatorname{CT}\!\left((1-x^2)P(x^2)^{m-k}P(x)^k\right).
\]
Every summand, after multiplication by $2/m$, vanishes modulo $4$,
since $v_2(\binom{m}{k}/m)\ge-v_2(k)$ and
$1+k-v_2(k)\ge2$. Differentiating $P^m$ gives
\[
 j c_{m,j}=m(c_{m-1,j-1}-c_{m-1,j+1}).
\]
Hence $2(M_{2m}-M_m)/m\equiv X_m\pmod4$.

Suppose $n\ge3$ is odd. The polynomial congruence
$P(x)^{2^n}\equiv P(x^h)^2\pmod4$ follows from the
Frobenius congruence modulo $2$ by one final squaring.
In the ring of formal Laurent series at zero, write $1/P(x)=\sum_j w_jx^j$.
One has $w_j=0$ for $j\le0$, and for positive $j$ its values in
residue classes $1,2,0$ modulo $3$ are $1,-1,0$, respectively,
since $1/P=x(1-x)/(1-x^3)$. The same period is valid at $j=0$.
Therefore, for $0\le j\le3$,
\[
 c_{m-1,j}\equiv c_{2a-1,j}
  -\sum_{s=1}^{j-1} c_{2a,s}w_{j-s}\pmod4.
\]
Indeed, the left-hand side is $\sum_s c_{2a,s}w_{j-hs}$ modulo $4$.
For $s\le0$ the weights agree with $w_{j-s}$ because
$h\equiv1\pmod3$. For $s\ge1$ they vanish since $h\ge4$.
Consequently
\[
\begin{split}
 X_m&\equiv X_{2a}-3c_{2a,1}+c_{2a,2}\\
    &=(1+a)L_{2a-1}-(2+6a)M_{2a-1}\pmod4 .
\end{split}
\]
For $a\equiv1\pmod4$ this reduces to $2L_{2a-1}$.
Modulo $2$, $P^{2a-1}=P\,P(x^2)^{a-1}$ implies
$L_{2a-1}\equiv c_{a-1,0}-c_{a-1,2}=M_{a-1}$.
As $4\mid a-1$, the coefficient $c_{a-1,2}$ vanishes modulo $2$
by Frobenius. Symmetry gives
$c_{a-1,0}\equiv P(1)^{a-1}\equiv1\pmod2$.
This proves the dyadic conclusion. The same computation gives
$v_2(M_{2m}-M_m)\ge n+1$ when $a\equiv3\pmod4$
and $n\ge3$ is odd.

\subsection{A ternary expansion at a double root}

Write $C_s=\binom{2s}{s}$, and define the integer
\[
 F_s=\binom{2s}{s}-\binom{2s}{s-2}
     = C_s\frac{2(2s+1)}{(s+1)(s+2)}
\]
with out-of-range binomial coefficients zero. The exact identity
$P=(1-x)^2/x+3$ gives
\[
 M_m=(-1)^m\sum_{k=0}^m(-3)^k\binom{m}{k}F_{m-k}.
\]
The first three terms are $(-1)^m C_m R(m)$, where
\[
 R(z)=\frac{2(2z+1)}{(z+1)(z+2)}-\frac{3z}{z+1}
       +\frac{9z(z-1)}{4(2z-1)} .
\]
Its expansion at zero belongs to $\mathbb Z_3[[z]]$ and begins
$R(z)=1-z/4+O(z^2)$.

Let $m=a3^n$, with $3\nmid a$. Legendre's formula gives
$v_3(C_m)=v_3(C_a)=V$. We prove
\begin{equation}\label{eq:ternary-lift}
 M_m\equiv(-1)^m C_m R(m)\pmod{3^{2n+V}}
 \qquad\bigl((3^n+1)/2\ge2n+V\bigr).
\end{equation}
For $k\ge(3^n+1)/2$, the factor $3^k$ proves this bound immediately.
If $3\le k<3^n/2$, write $v=v_3$. The falling-factorial ratio gives
\[
 v(C_{m-k}/C_m)
 =v\!\left(\frac{(m(m-1)\cdots(m-k+1))^2}
 {(2m)(2m-1)\cdots(2m-2k+1)}\right)
 =n-v(k)-v(C_k).
\]
Here $v(m-j)=v(2m-j)=v(j)$ for $0<j<3^n$.
Also $v\binom{m}{k}\ge n-v(k)$.
The denominators in $F_{m-k}$ cost $v(k-1)+v(k-2)$,
so the summand has valuation at least
\[
 2n+V+k-2v(k)-v(C_k)-v(k-1)-v(k-2).
\]
The last five terms together are nonnegative. For $k=3,4,5$ this
follows directly. For $k\ge6$ at most one of $k,k-1,k-2$ is
divisible by $3$, and Legendre's formula gives
$v(C_k)\le\lfloor\log_3(2k)\rfloor$. Thus the losses are
at most $3\lfloor\log_3(2k)\rfloor\le k$. For the last inequality
one checks $6\le k\le13$ directly. When $d=\lfloor\log_3(2k)\rfloor
\ge3$, use $k\ge3^d/2\ge3d$.
This proves \eqref{eq:ternary-lift}.

Since $(3^n+1)/2-2n$ increases for $n\ge1$, the threshold holds at all subsequent indices. Finally, cancelling the multiples of $3$ gives
\[
 \frac{C_{3m}}{C_m}
 =\prod_{\substack{1\le i\le3m\\3\nmid i}}(1+3m/i)
 \equiv1\pmod{3^{2n+2}}.
\]
Indeed, $\sum_{3\nmid i}1/i$ on this range belongs to
$3^{n+1}\mathbb Z_3$ by reflection $i\mapsto3m-i$,
and all terms of degree at least two in the product have the
required valuation. The threshold in \eqref{eq:ternary-lift}
persists at $n+1$, and $(-1)^{3m}=(-1)^m$.
Subtracting the two approximations gives
\[
 M_{3m}-M_m
 \equiv(-1)^m C_m(-m/2)\pmod{3^{2n+V}},
\]
because $R(3m)-R(m)=-m/2+O(m^2)$.
As $n\ge1$, the residue has valuation exactly $n+V$.
This proves Theorem~\ref{thm:lengyel55}.
\section{All ternary offsets}

The same analysis gives an exact order for every pair consisting of a
multiplier and an offset. In the following statement, $K_b$ is a
rational constant determined by the offset alone.

\begin{theorem}\label{thm:all-3-shifts}
Define $K_0=-1/4$, $K_1=7/12$, and, for $b\ge2$, set
\[
 (b+2)K_b=(2b+1)K_{b-1}+3(b-1)K_{b-2}
             +\frac{3(M_{b-1}+3M_{b-2})}{b+2}.
\]
For all positive $a$ coprime to $3$ and all $b\ge0$, there is
$n_0(a,b)$ such that every $n\ge n_0(a,b)$ satisfies
\[
 v_3(M_{a3^{n+1}+b}-M_{a3^n+b})
    =n+v_3\binom{2a}{a}+v_3(K_b).
\]
Moreover, $K_b>0$ for every $b\ge1$.
\end{theorem}

\begin{proof}
Put $m=a3^n$ and $V=v_3(C_a)$. We need a second starting value in
addition to \eqref{eq:ternary-lift}. In the corresponding expansion
of $M_{m+1}$, the terms with $0\le k\le3$ give
$(-1)^m C_m R_1(m)$, where
\[
 R_1(z)=-\frac{4(2z+1)(2z+3)}{(z+1)(z+2)(z+3)}
       +\frac{6(2z+1)}{z+2}-\frac{9z}{2}
       +\frac{9(z+1)z(z-1)}{4(2z-1)}.
\]
Thus $R_1(0)=1$ and $R'_1(0)=7/12$.

To bound the remaining terms $\binom{m+1}{k}(-3)^k F_{m+1-k}$,
suppose first that $4\le k\le(3^n+1)/2$
and write $\ell=k-1$, $v=v_3$. The same falling-factorial identity gives
$v(C_{m-\ell})=V+n-v(\ell)-v(C_\ell)$. The binomial identity
\[
 \binom{m+1}k=\frac{(m+1)m}{k(k-1)}\binom{m-1}{k-2}
\]
gives a second contribution of at least $n-v(k)-v(\ell)$.
The denominators in $F_{m-\ell}$ have valuation
$v(\ell-1)+v(\ell-2)$. The resulting summand has valuation at least
\[
 V+2n+k-v(k)-2v(k-1)-v(C_{k-1})-v(k-2)-v(k-3).
\]
The six final terms are bounded below independently of $k$, since
$v(j)\le\log_3 j$ and $v(C_j)\le\log_3(2j)$. For the remaining
$k>(3^n+1)/2$, the explicit $3^k$ gives valuation at least $2n+V$
for all sufficiently large $n$. Consequently there exists $d_1$
independent of $n$ such that
\[
 v_3(M_{m+1}-(-1)^m C_m R_1(m))\ge 2n+V-d_1
\]
for all sufficiently large $n$.

Put $R_0=R$ from \eqref{eq:ternary-lift} and define rational
functions regular at zero by
\[
 (z+b+2)R_b(z)=(2z+2b+1)R_{b-1}(z)+3(z+b-1)R_{b-2}(z)
 \qquad(b\ge2).
\]
The ordinary Motzkin recurrence has exactly these coefficients at $z=m$.
For each fixed $b$, the denominator $m+b+2$ eventually has the fixed
valuation $v_3(b+2)$. Induction therefore gives
\[
 v_3(M_{m+b}-(-1)^m C_m R_b(m))\ge2n+V-d_b
\]
for a constant $d_b$ and sufficiently large $n$.
Evaluation at zero gives $R_b(0)=M_b$. Differentiation gives the
stated recurrence for $K_b=R_b'(0)$, since the Motzkin recurrence
yields
\[
 2M_{b-1}+3M_{b-2}-M_b=\frac{3(M_{b-1}+3M_{b-2})}{b+2}.
\]
In particular $K_2=31/24$, so positivity follows from the recurrence
for all $b\ge1$. Hence every $K_b$ is nonzero.

For a fixed rational function regular at zero,
$v_3(R_b(3m)-R_b(m)-2mK_b)\ge2n-e_b$ for sufficiently large $n$,
with $e_b$ fixed. This follows by factoring $z^2$ after
subtracting the linear Taylor polynomial; the remaining rational
function is bounded on a neighborhood of zero. Using
$C_{3m}/C_m\equiv1\pmod{3^{2n+2}}$ from the previous proof,
the difference $M_{3m+b}-M_{m+b}$ equals
$(-1)^m C_m(2mK_b)$ up to an error of valuation at least
$2n+V-O_b(1)$. Its leading term has valuation $n+V+v_3(K_b)$,
which proves the theorem.
\end{proof}
\section{Provenance and scope}
The original graph question appears in Fajtlowicz's own catalogue, with
313 among the entries reported to have passed the small-graph tests
described after entry 107 \cite{fajt313}. Aouchiche and Hansen later
recorded experimental support \cite[Conj.~7.14]{ah313}.
The crown family occurs in their discussion of potential equality. A public
2 September 2026 computational note reports an exhaustive search on connected
triangle-free graphs of orders four through eleven, as well as unsuccessful searches
among larger graphs and perturbations of crowns \cite{village}.
Here it is the second factor of counterexamples. The use of a
Cartesian product for computing independence or distance spectra is
classical \cite{ikr313,indulal313,catral}. The contribution to compare
is the exact ratio controlled by $b(H)$, its application to the
original inequality, and its amplification on regular and Cayley graphs.

Lengyel's 2013 article states the full prime-power conjecture on p.~133.
Theorems~5.6 and 5.7 there give the multiplier-one results, also the subject
of his 2014 paper \cite{len13,len14}. We used the original theorem
statements and the follow-up article's published abstract to delimit
that prior result. For $p\ge5$ the original conjecture requires only
$a=1$. The present binary calculation covers $a\equiv1\pmod4$ at every
odd index in the source. The ternary calculation covers all $a$ with
$3\nmid a$ and extends to every fixed offset. The proof uses positivity
over the reals of the rational $K_b$ for $b\ge1$ only to establish
nonvanishing. The exact leading digit is determined in $\mathbb Q_3$.

The supplied corpus was inspected at the commit identified in
\cite{vmrc}. Manuscript [183] studies how a combinatorial inequality
changes under a graph operation and derives an explicit defect.
Manuscript [178] uses Frobenius digit separation for Laurent constant
terms, and [070] tracks exact valuation after a lift. These are
methodological starting points. The product inertia, normalized
coefficient identities and ternary noncancellation are derived here.

\begin{center}\small
\begin{tabularx}{\linewidth}{@{}p{.105\linewidth}X@{}}
\toprule
Test & Assessment\\ \midrule
G1, G2 & The external sources identify both numbered conjectures before this work.
The graph proof uses the original ordinary distance matrix, the original
independence number and the original triangle-free hypothesis. The
arithmetic proof retains the original index, primes and multiplier restrictions.\\
G3 & The 1988 catalogue entry and 2014 survey document the conjectural graph
bound. The 2013 arithmetic paper explicitly distinguishes the known
multiplier-one cases. Subsequent searches did not locate a full prior resolution
of either target. This is a bounded literature assessment.\\
G4 & The exact counting and spectral arguments prove the violation and
unbounded ratio. The binary coefficient identity and ternary remainder
bounds work over the full stated ranges. The Ramsey estimates use published
unconditional theorems.\\
G5 & The graph ratio has explicit growth on restricted classes. In the ternary
theorem a rigorously nonzero leading coefficient remains for every offset.
Both conclusions go beyond the separate input formulas.\\
G6 & The original lists, the distance-spectrum survey, product formulas,
public computational Graffiti reports, the corpus, Lengyel's bibliography
and searches by conjecture number and mathematical statement were compared.
No matching prior resolution or essentially identical proof was identified.
Historical priority and external review remain unverified.\\
\bottomrule
\end{tabularx}
\end{center}

\begin{thebibliography}{99}\small
\bibitem{fajt313} S. Fajtlowicz, \emph{Written on the Wall}, University of Houston, July 2004 compilation, pp.~80--81, Conjecture~313 in the block dated 27 August 1988. \url{https://independencenumber.wordpress.com/wp-content/uploads/2012/08/wow-july2004.pdf}.
\bibitem{ah313} M. Aouchiche and P. Hansen, Distance spectra of graphs: a survey, \emph{Linear Algebra Appl.} \textbf{458} (2014), 301--386, Conjecture~7.14. \url{https://doi.org/10.1016/j.laa.2014.06.010}.
\bibitem{ikr313} W. Imrich, S. Klav\v zar and D. F. Rall, \emph{Topics in Graph Theory: Graphs and Their Cartesian Product}, A K Peters, Wellesley, MA, 2008, Proposition~7.3.
\bibitem{indulal313} G. Indulal, Distance spectrum of graph compositions, \emph{Ars Math. Contemp.} \textbf{2} (2009), 93--100. \url{https://doi.org/10.26493/1855-3974.103.e09}.
\bibitem{erd313} P. Erd\H{o}s, Graph theory and probability. II, \emph{Canad. J. Math.} \textbf{13} (1961), 346--352. \url{https://doi.org/10.4153/CJM-1961-029-9}.

\bibitem{village} AI Village, \emph{WOW-I 235, 237, 313: investigated and declined}, computational note dated 2 September 2026, section on 313. \href{https://gitlab.com/ai-village-agents/village/graffiti-verification/-/blob/8bd9f691fd56a2ba070ba3a2cb3b6c8750ad6d71/verify/notes_wow1_235_237_313_2026-09-02.md}{Versioned note, commit 8bd9f69}.
\bibitem{catral} M.~Catral, L.~Ciardo, L.~Hogben and C.~Reinhart, Spectra of products of digraphs, \emph{Electron. J. Linear Algebra} \textbf{36} (2020). \href{https://doi.org/10.13001/ela.2020.5243}{doi:10.13001/ela.2020.5243}.
\bibitem{abs} N.~Alon, S.~Ben-Shimon and M.~Krivelevich, A note on regular Ramsey graphs, \emph{J. Graph Theory} \textbf{64} (2010), 244--249. \href{https://doi.org/10.1002/jgt.20453}{doi:10.1002/jgt.20453}.
\bibitem{alon} N.~Alon, Lov\'asz, vectors, graphs and codes, Theorem~2.1. \href{https://web.math.princeton.edu/~nalon/PDFS/ll70.pdf}{Author's copy with the explicit connection set}.
\bibitem{len13} T.~Lengyel, On divisibility properties of some differences of Motzkin numbers, \emph{Ann. Math. Inform.} \textbf{41} (2013), 121--136. \href{https://sites.oxy.edu/lengyel/papers/AMI_41_from121to136.pdf}{Author's copy}.
\bibitem{len14} T.~Lengyel, Exact $p$-adic orders for differences of Motzkin numbers, \emph{Int. J. Number Theory} \textbf{10} (2014), 653--667. \href{https://doi.org/10.1142/S1793042113501157}{doi:10.1142/S1793042113501157}.
\bibitem{vmrc} QEEDDD, \emph{Verifiable Mathematical Reasoning Corpus}, manuscripts [070], [178] and [183], \href{https://github.com/QEEDDD/Verifiable-Mathematical-Reasoning-Corpus/tree/5887cdea4898f0a3aa5bf7a8e0829697bd2b0f4b}{commit 5887cdea4898f0a3aa5bf7a8e0829697bd2b0f4b}.

\end{thebibliography}
\end{document}
